\section{Conclusions}\label{sec:conclusions}

We set out to measure, in one engine and over \nCells{} cells of \nScenes{} scenes from one million to one billion points, where the query speed of a point-cloud spatial index comes from. The answer has three parts. The primitive and its leaf account for the whole gain between a library default and the best structure a search can find, a factor of \leverOctreeDefaultTuned{} over an octree default and \leverKdDefaultTuned{} over a median kd-tree; a single split on the scene's flatness picks the primitive within \ruleRegret{} to \ruleRegretTimed{} of a per-cloud search, and nanoflann's default already embodies the right choice for three-dimensional scenes. Composing primitives into nested hierarchies adds nothing once the hierarchy is in RAM: \nestedPickedAll{} of \nCells{} cells in memory, none of seven workloads, none on the GPU, none under emulated node costs, and \oocInRamAdopted{} of \oocInRamCells{} out-of-core cells with the leaves on a device. It pays in \oocPagedAdopted{} of \oocPagedCells{} cells with the hierarchy itself paged from a slow device, which is the regime of the result it revisits, and a regime that pinning a few tens of megabytes beats several times over. The mechanism is that a node visit costs an order of magnitude more at a leaf than near the root, so the visits a wide prefix saves are the cheap ones, and node counts do not predict time; the memory layout does, and two-megabyte pages buy more than a second primitive ever did.

The study has limits of its own. It is one engine on two machines; its devices are modelled; its traces come from three pipeline stages; and its rule is fitted on sixteen scenes. What it offers in return is a controlled protocol, a released instrument with its traces, configurations and scripts, and a message that is smaller than the one it replaces and easier to act on: choose the primitive from the geometry, the leaf from the device, the layout from the page size, and measure the rest.
